\documentclass[11pt]{article}

\usepackage[margin=1in]{geometry}
\usepackage{amsmath}
\usepackage{amssymb}
\usepackage{amsthm}
\usepackage{booktabs}
\usepackage{array}
\usepackage{enumitem}
\usepackage{hyperref}
\hypersetup{colorlinks=true, linkcolor=blue, citecolor=blue, urlcolor=blue}

\newtheorem{definition}{Definition}[section]
\newtheorem{example}[definition]{Example}
\newtheorem{proposition}[definition]{Proposition}
\newtheorem{lemma}[definition]{Lemma}
\newtheorem{remark}[definition]{Remark}
\newtheorem{corollary}[definition]{Corollary}

\newcommand{\N}{N}
\newcommand{\Mset}{M}

\title{Pareto Optimality Is Incompatible with a Bounded Per-Agent Subsidy:\\
       A Minimal Example}
\author{Mainak Sarkar}
\date{\today}

\begin{document}
\maketitle

\begin{abstract}
We exhibit a two-agent, three-chore instance with binary marginal costs in
which \emph{no} complete allocation is simultaneously (i) envy-free once each
agent is paid a subsidy of at most one unit, and (ii) Pareto optimal. The two
properties are not merely hard to achieve together on this instance --- they
are mutually exclusive: the set of allocations satisfying the subsidy
guarantee and the set of Pareto optimal allocations are disjoint. We use the
standard definition of Pareto optimality of Chakraborty, Igarashi, Suksompong
and Zick~\cite{CISZ21}, translated from goods to chores by the usual sign
flip, and we spell out one point where the two settings genuinely differ ---
allocations must be taken complete for the chores version of the definition
to be non-vacuous.
\end{abstract}

\section{Summary, for the impatient reader}
\label{sec:summary}

\begin{itemize}[leftmargin=1.4em]
  \item Two agents share the cost function $c(S) = \min(|S|,2)$ on three
        chores: the first two chores hurt, the third is free once you already
        hold two.
  \item Concentrating all three chores on one agent is the \emph{cheapest}
        outcome in total (cost $2$), because the third chore is absorbed for
        free. Splitting the chores always costs $3$ in total.
  \item But concentration creates an envy gap of exactly $2$: the loaded
        agent needs a subsidy of $2$ to stop envying the agent holding
        nothing. A guarantee capping the subsidy at $1$ per agent forbids
        this.
  \item Every allocation that respects the cap is a split, and every split is
        Pareto dominated by concentrating the chores.
  \item Conclusion: on this instance, \textbf{no allocation is both Pareto
        optimal and feasible under the $\{0,1\}$-subsidy guarantee}. This is
        not an artifact of a bad tie-break; the two requirements are
        logically incompatible here.
\end{itemize}

\section{The setting}
\label{sec:setting}

There are $\N = \{1,\dots,n\}$ agents and $\Mset$ a finite set of indivisible
\emph{chores}. Each agent $i$ has a cost function $c_i : 2^{\Mset} \to
\mathbb{Z}_{\ge 0}$ with $c_i(\emptyset) = 0$. We say $c_i$ has \emph{binary
marginals} if
\[
  c_i(S \cup \{e\}) - c_i(S) \in \{0,1\}
  \qquad \text{for every } S \subseteq \Mset \text{ and every } e \notin S.
\]
No further structure is assumed: $c_i$ need not be additive, submodular, or
subadditive, and different agents' cost functions need not be related.

A \emph{complete allocation} $A = (A_1,\dots,A_n)$ partitions $\Mset$ among
the agents. A subsidy vector $p \in \mathbb{R}^n_{\ge 0}$ makes $(A,p)$
\emph{envy-free} if
\[
  c_i(A_i) - p_i \;\le\; c_i(A_j) - p_j \qquad \text{for all } i,j \in \N .
\]

We are interested in the following guarantee, which caps the subsidy paid to
any single agent:
\begin{definition}[The $\{0,1\}$-subsidy property]
\label{def:S}
A complete allocation $A$ is \emph{$S$-feasible} if there exists
$p \in \{0,1\}^n$ making $(A,p)$ envy-free.
\end{definition}
Definition~\ref{def:S} is the property that a subsidy of at most one unit per
agent suffices to eliminate all envy under $A$. This note asks whether an
$S$-feasible allocation can always be chosen to be Pareto optimal as well.

\section{Pareto optimality}
\label{sec:po}

\subsection{The standard definition, for goods}

We use the definition of Chakraborty, Igarashi, Suksompong and
Zick~\cite[Sec.~2]{CISZ21}, stated for their setting of \emph{goods}, where
each agent has a value function $v_i : 2^{\Mset} \to \mathbb{R}_{\ge 0}$
(higher is better):

\begin{quote}
\emph{An allocation $A'$ is said to Pareto dominate an allocation $A$ if
$v_i(A_i') \ge v_i(A_i)$ for all agents $i \in \N$ and $v_j(A_j') > v_j(A_j)$
for some agent $j \in \N$. An allocation is Pareto optimal (or PO for short)
if it is not Pareto dominated by any other allocation.}
\end{quote}

\subsection{The chores version}

Our setting is \emph{chores}, so we work with costs rather than values. The
standard translation between the two is $v_i = -c_i$: a chore is something an
agent does not want, so cost is disutility. Substituting $v_i = -c_i$ into the
inequalities above,
\[
  v_i(A_i') \ge v_i(A_i)
  \iff -c_i(A_i') \ge -c_i(A_i)
  \iff c_i(A_i') \le c_i(A_i),
\]
and the strict inequality flips the same way. The chores form of the
definition is therefore \emph{the same statement}, not an independent one:

\begin{definition}[Pareto domination, chores form]
\label{def:podom}
Allocation $B$ \emph{Pareto dominates} allocation $A$ if
$c_i(B_i) \le c_i(A_i)$ for every $i \in \N$, with strict inequality for at
least one $j \in \N$. An allocation is \emph{Pareto optimal (PO)} if no
allocation Pareto dominates it.
\end{definition}

\subsection{Why allocations must be taken complete}
\label{sec:completeness}

Chakraborty et al.\ allow \emph{incomplete} allocations --- some items may go
unassigned --- and Pareto optimality is defined over all allocations,
complete or not. For goods this costs nothing: withholding a good from every
agent can only make everyone worse off (assuming positive value), so an
incomplete allocation is automatically dominated by completing it, and hence
can never be PO. Chakraborty et al.\ note exactly this in their proof of
Theorem~4.2.

For chores the same argument runs backward. Withholding a chore makes
\emph{every} agent weakly better off. If incomplete allocations were allowed
to compete as dominators, \emph{every} complete allocation would be dominated
by the empty allocation (assign nothing to anybody), and the notion of Pareto
optimality would collapse to nothing being PO except the empty allocation
itself --- a vacuous statement that says nothing about how to divide the
chores that must, in fact, be divided.

We therefore adopt the standard convention of the chore fair-division
literature and restrict Definition~\ref{def:podom} to \emph{complete}
allocations throughout: an allocation is a partition of $\Mset$, and
domination is checked only against other partitions of $\Mset$. This is not a
weaker or different notion of Pareto optimality; it is the same definition,
applied to the domain on which it is non-vacuous for chores.

\section{The instance}
\label{sec:instance}

\begin{example}
\label{ex:main}
Let $n = 2$ agents and $\Mset = \{a,b,c\}$, three chores. Both agents share
the same cost function
\[
  c(S) \;=\; \min(|S|,\,2).
\]
\end{example}

\begin{table}[h]
\centering
\begin{tabular}{@{}lcccc@{}}
\toprule
$|S|$ & $0$ & $1$ & $2$ & $3$ \\
\midrule
$c(S)$ & $0$ & $1$ & $2$ & $2$ \\
\bottomrule
\end{tabular}
\caption{The shared cost function. The first two chores cost one unit each;
the third is free once the agent already holds two.}
\label{tab:cost}
\end{table}

\begin{proposition}
\label{prop:inclass}
The cost function of Example~\ref{ex:main} has binary marginals, is
monotone, satisfies $c(\emptyset)=0$, and is \emph{not} additive.
\end{proposition}

\begin{proof}
$c(\emptyset) = \min(0,2) = 0$. Monotonicity is immediate from the definition
of $\min$. For the marginals: adding a chore to the empty bundle gives
$c(\{e\}) - c(\emptyset) = 1 - 0 = 1$; adding a chore to a singleton gives
$c(S) - c(\{e'\}) = 2 - 1 = 1$; adding a chore to a bundle of size two gives
$c(\Mset) - c(S) = 2 - 2 = 0$. Every marginal is in $\{0,1\}$.

If $c$ were additive, we would need $c(\{a,b,c\}) = c(\{a\}) + c(\{b\}) +
c(\{c\}) = 1+1+1 = 3$. But $c(\{a,b,c\}) = \min(3,2) = 2 \ne 3$, so $c$ is not
additive.
\end{proof}

\section{The eight allocations}
\label{sec:eight}

With $n=2$ and $|\Mset|=3$, there are $2^3 = 8$ complete allocations. Because
both agents share the same cost function, the envy-free-with-subsidy
condition simplifies sharply.

\begin{lemma}[Forced subsidy gap]
\label{lem:gap}
For $n=2$ agents with a shared cost function $c$, an allocation $A$ admits an
envy-free subsidy vector $p$ if and only if
\[
  p_1 - p_2 \;=\; c(A_1) - c(A_2).
\]
\end{lemma}

\begin{proof}
Envy-freeness requires both $c(A_1) - p_1 \le c(A_2) - p_2$ (agent $1$ does
not envy agent $2$) and $c(A_2) - p_2 \le c(A_1) - p_1$ (agent $2$ does not
envy agent $1$). Rearranging, the first gives $p_1 - p_2 \ge c(A_1) - c(A_2)$
and the second gives $p_1 - p_2 \le c(A_1) - c(A_2)$. Together these force
equality.
\end{proof}

Lemma~\ref{lem:gap} leaves no freedom at all: once $A$ is fixed, the
\emph{difference} $p_1 - p_2$ is determined exactly. Since $p \in
\{0,1\}^2$, the achievable differences are only $\{-1,0,1\}$; the allocation
is $S$-feasible (Definition~\ref{def:S}) exactly when $|c(A_1)-c(A_2)| \le 1$.

\begin{table}[h]
\centering
\begin{tabular}{@{}llccccl@{}}
\toprule
$A_1$ & $A_2$ & $(c(A_1),c(A_2))$ & $p_1-p_2$ forced & $S$-feasible? & PO? & dominated by \\
\midrule
$\emptyset$   & $\{a,b,c\}$ & $(0,2)$ & $-2$ & \textbf{no} & \textbf{yes} & --- \\
$\{a\}$       & $\{b,c\}$   & $(1,2)$ & $-1$ & yes & no & $(\emptyset,\{a,b,c\})$ \\
$\{b\}$       & $\{a,c\}$   & $(1,2)$ & $-1$ & yes & no & $(\emptyset,\{a,b,c\})$ \\
$\{c\}$       & $\{a,b\}$   & $(1,2)$ & $-1$ & yes & no & $(\emptyset,\{a,b,c\})$ \\
$\{a,b\}$     & $\{c\}$     & $(2,1)$ & $1$  & yes & no & $(\{a,b,c\},\emptyset)$ \\
$\{a,c\}$     & $\{b\}$     & $(2,1)$ & $1$  & yes & no & $(\{a,b,c\},\emptyset)$ \\
$\{b,c\}$     & $\{a\}$     & $(2,1)$ & $1$  & yes & no & $(\{a,b,c\},\emptyset)$ \\
$\{a,b,c\}$   & $\emptyset$ & $(2,0)$ & $2$  & \textbf{no} & \textbf{yes} & --- \\
\bottomrule
\end{tabular}
\caption{All eight complete allocations of Example~\ref{ex:main}. Exactly two
are Pareto optimal (the concentrated ones); exactly six are $S$-feasible (the
splits). The two sets do not overlap.}
\label{tab:eight}
\end{table}

\section{The two Pareto optimal allocations}
\label{sec:po-two}

\begin{proposition}
\label{prop:po}
The allocations $(\{a,b,c\},\emptyset)$ and $(\emptyset,\{a,b,c\})$ are
Pareto optimal, and no other allocation is.
\end{proposition}

\begin{proof}
Let $A = (\{a,b,c\},\emptyset)$, with costs $(2,0)$. Suppose $B$ dominates
$A$: then $c(B_1) \le 2$ and $c(B_2) \le 0$, with at least one strict. Costs
are non-negative, so $c(B_2) \le 0$ forces $c(B_2) = 0$, and the only subset
with cost $0$ is $\emptyset$ (Table~\ref{tab:cost}), so $B_2 = \emptyset$.
Completeness then forces $B_1 = \{a,b,c\}$. So $B = A$: no allocation other
than $A$ itself satisfies the domination inequalities, and $A$ is PO. The
argument for $(\emptyset,\{a,b,c\})$ is symmetric.

For every other allocation, Table~\ref{tab:eight} exhibits an explicit
dominator. Take $(\{a,b\},\{c\})$, with costs $(2,1)$, as a representative
case: $c(\{a,b\}) = 2 = c(\{a,b,c\})$, while $c(\{c\}) = 1 > 0 = c(\emptyset)$.
So $(\{a,b,c\},\emptyset)$ dominates it --- agent $1$ is unaffected because her
bundle is already saturated, and agent $2$ strictly improves. The same pattern
verifies every other row. Hence no allocation outside
$\{(\{a,b,c\},\emptyset), (\emptyset,\{a,b,c\})\}$ is PO.
\end{proof}

\begin{remark}[The two PO allocations do not dominate each other]
\label{rem:incomparable}
It is worth checking directly that $(\{a,b,c\},\emptyset)$, with costs
$(2,0)$, does \emph{not} Pareto dominate $(\emptyset,\{a,b,c\})$, with costs
$(0,2)$ --- and vice versa. Domination of $(0,2)$ by $(2,0)$ would require
$2 \le 0$ (agent $1$'s cost) \emph{and} $0 \le 2$ (agent $2$'s cost); the
first fails. Domination in the other direction would require $0 \le 2$
\emph{and} $2 \le 0$; again the second fails. So neither dominates the other.
This is not a contradiction: Pareto optimality is a partial order, not a
ranking, and $(2,0)$ and $(0,2)$ are simply two incomparable points on the
Pareto frontier --- moving from one to the other helps one agent while
hurting the other, which is a trade-off, not an improvement. Both are
undominated by anything, including each other, which is exactly what makes
both of them Pareto optimal simultaneously.
\end{remark}

\section{The incompatibility}
\label{sec:incompat}

\begin{corollary}[Main observation]
\label{cor:main}
On the instance of Example~\ref{ex:main}, no complete allocation is both
$S$-feasible (Definition~\ref{def:S}) and Pareto optimal
(Definition~\ref{def:podom}).
\end{corollary}

\begin{proof}
By Proposition~\ref{prop:po}, the Pareto optimal allocations are exactly
$(\{a,b,c\},\emptyset)$ and $(\emptyset,\{a,b,c\})$, with forced subsidy gaps
$p_1-p_2 = \pm 2$ (Lemma~\ref{lem:gap} and Table~\ref{tab:eight}). Neither gap
is achievable by any $p \in \{0,1\}^2$, whose achievable gaps are
$\{-1,0,1\}$. So neither PO allocation is $S$-feasible. Every other
allocation is $S$-feasible (Table~\ref{tab:eight}) but, by
Proposition~\ref{prop:po}, is not PO. Hence the two sets --- $S$-feasible
allocations and Pareto optimal allocations --- are disjoint.
\end{proof}

\begin{remark}[Why: efficiency and fairness pull in opposite directions]
The two facts behind Corollary~\ref{cor:main} are two sides of the same
feature of $c$. Because $c$ saturates at $|S|=2$, the third chore is
\emph{free} to an agent who already holds two: moving it onto that agent
costs nothing and relieves whoever else was holding it, which is exactly what
makes concentration efficient (Pareto optimal). But that same saturation
means the loaded agent's cost jumps straight from $0$ to $2$ relative to the
empty agent, an envy gap the $\{0,1\}$-subsidy guarantee cannot fund. The
guarantee therefore forces a split; every split leaves one chore un-absorbed
that someone else would have taken for free, which is exactly the source of
the domination in Table~\ref{tab:eight}.
\end{remark}

\begin{remark}[The efficiency cost of the guarantee]
The cheapest Pareto optimal allocation has total cost $c(A_1)+c(A_2) = 2$.
The cheapest $S$-feasible allocation has total cost $3$ (every split costs
$2+1$ in some order). Enforcing the $\{0,1\}$-subsidy guarantee therefore
costs society one extra unit of total disutility on this instance --- a
$50\%$ overhead relative to the efficient outcome.
\end{remark}

\section{Further observations (informal)}
\label{sec:further}

The following two points were checked computationally and are recorded here
for context; they are not needed for Corollary~\ref{cor:main}, which is
already fully self-contained on the three-chore instance above.

\begin{itemize}[leftmargin=1.4em]
  \item \textbf{The phenomenon is not an artifact of $n=2,m=3$.} The same
        construction, $c(S) = \min(|S|,t)$ with all agents sharing $c$,
        produces analogous instances --- no $S$-feasible allocation is
        Pareto optimal --- for a range of $(n,m,t)$ with $t=2$ and
        $m \ge n+1$, checked exhaustively for $n \in \{2,3,4\}$.
  \item \textbf{Additivity seems to be exactly what is missing.} In every
        instance checked (exhaustive search at small sizes, and random
        sampling at larger sizes) where every agent's cost function is
        additive, an $S$-feasible allocation that is also Pareto optimal was
        always found. Proposition~\ref{prop:inclass} confirms that the
        instance above is emphatically not additive; that non-additivity
        appears to be essential to the incompatibility.
\end{itemize}

\begin{thebibliography}{9}
\bibitem{CISZ21}
Mithun Chakraborty, Ayumi Igarashi, Warut Suksompong, and Yair Zick.
\emph{Weighted Envy-Freeness in Indivisible Item Allocation}.
ACM Transactions on Economics and Computation, 9(3):1--18, 2021.
\end{thebibliography}

\end{document}
